% Shared body of the 250-word abstract (manuscript and resumen_250_palabras.tex).
La evapotranspiración real ($\ETa$) estimada con modelos de balance
de energía como METRIC suele forzarse con una sola estación meteorológica,
supuesto débil en territorios heterogéneos. Se evaluó si regionalizar ese
forzamiento con regresión geográficamente ponderada (GWR) modifica la $\ETa$
frente a una estación representativa (C1), una regresión global y la
interpolación por distancia inversa (IDW). Se usó EarthFlux, un paquete de R
de código abierto, con una escena MODIS de la provincia de Cauquenes (Chile,
15-02-2019), \NStations{} estaciones públicas y validación cruzada espacial
anidada en \NCVDays{} días. GWR superó a C1 y a la regresión global, pero no a
IDW: el RMSE de la temperatura del aire fue \CVTaCOneRmse{}, \CVTaCTwoOLSRmse{},
\CVTaCThreeGWRRmse{} y \CVTaCTwoIDWRmse~\si{\celsius}, respectivamente. GWR requirió
más estaciones, extrapoló temperaturas fuera del rango observado en el
\RangeTaCThreeGWRAbove\,\% de las celdas y no fue factible para el viento.
En METRIC, con las mismas anclas, GWR redujo la $\ETa$ media en
\DEtaCThreeGWRPct\,\% respecto de C1. La calibración interna absorbió el
efecto sobre el balance instantáneo ($\ETrF$: \DEtrfCThreeGWRPct\,\%), y la
diferencia entró por la $\ETr$ diaria. Ese efecto superó \RatioSA{} veces la
incertidumbre por calibración y \RatioSB{} veces la debida al ruido MODIS.
Sin mediciones independientes de flujos, la exactitud de $\ETa$ no pudo
evaluarse. El forzamiento meteorológico es una fuente de incertidumbre de
primer orden que debe declararse en toda estimación de $\ETa$ usada para
gestionar el agua; la mayor complejidad de GWR no garantizó mayor exactitud.

\medskip\noindent\textbf{Palabras clave:} evapotranspiración real; METRIC;
MODIS; regresión geográficamente ponderada; validación cruzada espacial;
software científico abierto; Chile central.
